% GENERATED FILE. Edit canonical lessons, then run npm run generate:books.
% canonical-source-hash: fnv1a64:437f025ba6edfe24
% canonical-lessons: FR-C60-pleurer, FR-C60-laver, FR-C60-cuisiner, FR-C60-jouer, FR-C60-apprendre

\chapter{Cry, Wash, Cook, Play, Learn}
\label{ch:fr-cry-wash-cook-play-learn}
\hlchaptermodality{60}

\begin{chapteropening}
\textbf{By the end of this chapter:} \emph{I can say to cry, to wash, to cook, to play, to learn.}

Complete the last lesson of chapter 60: I can say to cry, to wash, to cook, to play, to learn.
\end{chapteropening}

\section[pleurer]{pleurer — to cry}

\label{lesson:FR-C60-pleurer}



\begin{quote}
Before the new one: say the French for a teacher, then the French for to laugh.
\end{quote}

\subsection*{You'll want to know: pleurer}
\textbf{pleurer} — \textquotedblleft{}to cry\textquotedblright{}.

\begin{cousinweb}
\textbf{pleurer} $\leftarrow$ Old French \emph{plorer} $\leftarrow$ Latin \emph{plorare}, \textquotedblleft{}to wail, to weep aloud\textquotedblright{} — inherited. (EtymDB 2.1)
\end{cousinweb}

\begin{grammarlens}[title={What you've built}]
One more word for this chapter.
\end{grammarlens}

\subsection*{Your turn}
\begin{itemize}
\raggedright
  \item \emph{Say it:} \emph{pleurer}
  \item \emph{Say it:} \emph{pleurer}, once more
  \item \emph{Say it:} say \emph{pleurer}
  \item \emph{Recall:} say the French for a teacher, then the French for to laugh, then say \emph{pleurer} again
\end{itemize}

\begin{tcolorbox}[breakable,colback=teal!4,colframe=teal!35!black,title={Before you move on}]
What does pleurer mean? (\textquotedblleft{}To cry\textquotedblright{}.) Say it once more.
\end{tcolorbox}
\section[laver]{laver — to wash}

\label{lesson:FR-C60-laver}



\begin{quote}
Before the new one: say the French for to laugh, then the French for to cry.
\end{quote}

\subsection*{You'll want to know: laver}
\textbf{laver} — \textquotedblleft{}to wash\textquotedblright{}.

\begin{cousinweb}
\textbf{laver} $\leftarrow$ Latin \emph{lavare}, \textquotedblleft{}to wash\textquotedblright{} — inherited. (EtymDB 2.1)
\end{cousinweb}

\begin{grammarlens}[title={What you've built}]
One more word for this chapter.
\end{grammarlens}

\subsection*{Your turn}
\begin{itemize}
\raggedright
  \item \emph{Say it:} \emph{laver}
  \item \emph{Say it:} \emph{laver}, once more
  \item \emph{Say it:} say \emph{laver}
  \item \emph{Recall:} say the French for to laugh, then the French for to cry, then say \emph{laver} again
\end{itemize}

\begin{tcolorbox}[breakable,colback=teal!4,colframe=teal!35!black,title={Before you move on}]
What does laver mean? (\textquotedblleft{}To wash\textquotedblright{}.) Say it once more.
\end{tcolorbox}
\section[cuisiner]{cuisiner — to cook}

\label{lesson:FR-C60-cuisiner}



\begin{quote}
Before the new one: say the French for to cry, then the French for to wash.
\end{quote}

\subsection*{You'll want to know: cuisiner}
\textbf{cuisiner} — \textquotedblleft{}to cook\textquotedblright{}.

\begin{cousinweb}
\textbf{cuisiner} — built from the French word for \textquotedblleft{}kitchen, cooking\textquotedblright{}, which $\leftarrow$ Vulgar Latin \emph{*cocina} $\leftarrow$ Latin \emph{coquina}, \textquotedblleft{}kitchen\textquotedblright{}, from \emph{coquere}, \textquotedblleft{}to cook\textquotedblright{}. (EtymDB 2.1)
\end{cousinweb}

\begin{grammarlens}[title={What you've built}]
One more word for this chapter.
\end{grammarlens}

\subsection*{Your turn}
\begin{itemize}
\raggedright
  \item \emph{Say it:} \emph{cuisiner}
  \item \emph{Say it:} \emph{cuisiner}, once more
  \item \emph{Say it:} say \emph{cuisiner}
  \item \emph{Recall:} say the French for to cry, then the French for to wash, then say \emph{cuisiner} again
\end{itemize}

\begin{tcolorbox}[breakable,colback=teal!4,colframe=teal!35!black,title={Before you move on}]
What does cuisiner mean? (\textquotedblleft{}To cook\textquotedblright{}.) Say it once more.
\end{tcolorbox}
\section[jouer]{jouer — to play}

\label{lesson:FR-C60-jouer}



\begin{quote}
Before the new one: say the French for to wash, then the French for to cook.
\end{quote}

\subsection*{You'll want to know: jouer}
\textbf{jouer} — \textquotedblleft{}to play\textquotedblright{}.

\begin{cousinweb}
\textbf{jouer} $\leftarrow$ Old French \emph{joer} $\leftarrow$ Latin \emph{iocari}, \textquotedblleft{}to jest, to joke\textquotedblright{} — inherited. (EtymDB 2.1)
\end{cousinweb}

\begin{grammarlens}[title={What you've built}]
One more word for this chapter.
\end{grammarlens}

\subsection*{Your turn}
\begin{itemize}
\raggedright
  \item \emph{Say it:} \emph{jouer}
  \item \emph{Say it:} \emph{jouer}, once more
  \item \emph{Say it:} say \emph{jouer}
  \item \emph{Recall:} say the French for to wash, then the French for to cook, then say \emph{jouer} again
\end{itemize}

\begin{tcolorbox}[breakable,colback=teal!4,colframe=teal!35!black,title={Before you move on}]
What does jouer mean? (\textquotedblleft{}To play\textquotedblright{}.) Say it once more.
\end{tcolorbox}
\section[apprendre]{apprendre — to learn}

\label{lesson:FR-C60-apprendre}



\begin{quote}
Before the new one: say the French for to cook, then the French for to play.
\end{quote}

\subsection*{You'll want to know: apprendre}
\textbf{apprendre} — \textquotedblleft{}to learn\textquotedblright{}.

\begin{cousinweb}
\textbf{apprendre} $\leftarrow$ Old French \emph{aprendre} $\leftarrow$ Latin \emph{apprehendere}, \textquotedblleft{}to seize, to grasp\textquotedblright{} — inherited: to learn is to grasp. (EtymDB 2.1) It is \emph{prendre}, \textquotedblleft{}to take\textquotedblright{}, with \emph{a-} in front.
\end{cousinweb}

\begin{grammarlens}[title={What you've built}]
One more word for this chapter.
\end{grammarlens}

\subsection*{Your turn}
\begin{itemize}
\raggedright
  \item \emph{Say it:} \emph{apprendre}
  \item \emph{Say it:} \emph{apprendre}, once more
  \item \emph{Say it:} say \emph{apprendre}
  \item \emph{Recall:} say the French for to cook, then the French for to play, then say \emph{apprendre} again
\end{itemize}

\begin{tcolorbox}[breakable,colback=teal!4,colframe=teal!35!black,title={Before you move on}]
What does apprendre mean? (\textquotedblleft{}To learn\textquotedblright{}.) Say it once more.
\end{tcolorbox}
